%%%%%%%%%%%%%%%%%%%%%%%%%%%%%%%%%%%%%%%%%%%%%%%%%%%%%%%%%%%%%%%%%%%%%%%%%%%%%%%%%%%%%%%%%%%%%%%%%
% Chapter 7: Conclusion   (the \chapter command is in main.tex)
%%%%%%%%%%%%%%%%%%%%%%%%%%%%%%%%%%%%%%%%%%%%%%%%%%%%%%%%%%%%%%%%%%%%%%%%%%%%%%%%%%%%%%%%%%%%%%%%%

\label{sec:c_objectives}This project let a delivery drone decide when to fly home using an Extended Kalman Filter that estimates the battery's state of charge and how uncertain it is. All four objectives were met. The MATLAB simulator reproduced realistic battery behaviour, including a voltage sag of 0.37\,V per cell at take-off, and its battery deliberately differs from the filter's model. Learning the internal resistance during flight made the EKF accurate and honest: on a normal flight its SoC error was 0.97\% against 4.36\% for charge counting, and it corrected a 15\% wrong start within 42\,s while charge counting stayed wrong all flight. The full filter, which also learns the sensor bias and capacity, was the only one honest in every normal scenario, including an aged battery on which the most accurate filter overestimated the charge by 8.5\%. The risk-aware return-to-home rule built on it brought the drone home in all 500 random missions and broke the 10\% reserve in only 2.8\% of them, within its 5\% promise, while charge counting let 30\% of drones run out and the same filter used without its uncertainty broke the reserve in 38\% of missions. At every threshold tested the rule kept its promise, so its probabilities can be trusted.

\label{sec:c_next}The work is limited to simulation and a representative battery curve; the most useful next steps are calibrating the sensor bias before take-off, estimating the capacity between flights and testing on real flight data (Section~\ref{sec:d_limitations}).

The main lesson is that, for a safety decision, a filter's uncertainty is as valuable as its estimate, provided that uncertainty is honest.
